\documentclass[11pt]{article}
\usepackage[margin=0.9in]{geometry}
\usepackage{amsmath,amssymb,amsthm,booktabs,mathrsfs}
\usepackage[T1]{fontenc}
\usepackage{lmodern,microtype}
\usepackage[colorlinks=true,linkcolor=blue,urlcolor=blue]{hyperref}
\newtheorem{theorem}{Theorem}[section]
\newtheorem{lemma}[theorem]{Lemma}
\newtheorem{proposition}[theorem]{Proposition}
\newtheorem{corollary}[theorem]{Corollary}
\theoremstyle{remark}
\newtheorem{remark}[theorem]{Remark}
\newcommand{\norm}[1]{\left\|#1\right\|}
\title{Toward a four-dimensional Yang--Mills construction\\
\large Notebook 18: integrating overlapping cutoff interactions}
\author{Research working document}
\date{17 October 2026}
\setlength{\emergencystretch}{5em}
\begin{document}
\maketitle
\begin{abstract}
We extend the Gaussian cutoff expansion to interactions whose
supports overlap and can have arbitrarily many cells. An explicit
weighted local norm controls their expansion. The covariance forest
is applied after grouping overlapping factors, but its derivatives
retain their original cell endpoints. This prevents uncontrolled
powers of the grouped support size at high-degree vertices. We
obtain a nonvanishing partition function, a holomorphic logarithm
with summable activities, and a bound linear in the interaction norm
for its difference from the pure Gaussian cutoff logarithm.
A support-algebra construction also gives strictly local background
activities when the precision has a small local decomposition.
For a real baseline, a tunable relative Gaussian comparison reduces
the density cost per cell from a fixed exponential in its dimension
to an exponential with an arbitrarily small coefficient. The
theorems specify the covariance and coefficient estimates that
the Yang--Mills quantum step must supply, and a final section
identifies the corresponding original Ba\l aban integration route.
\end{abstract}

\section{The integration problem addressed here}
\label{oc18-target}

The nonlinear Wilson action and its coordinate factors generate
interactions supported on several overlapping cells. Treating each
connected collection as a single large Gaussian cell can introduce
a loss growing with both its support size and its branching degree.
The proof below keeps the original cell structure through the
covariance derivatives and sums the resulting degree factorials
exactly. The original Gaussian cell dimension stays fixed.

This is an abstract integration theorem with explicit hypotheses.
The homogeneous flat pivot Gaussian has the spectral and locality
estimates from Notebooks 6 and 7; Notebook 17 records its joint
fluctuation bounds. To apply the present theorem to the complete
Yang--Mills integrand, the action, Haar and Jacobian contributions
must be represented with the stated weighted local coefficient
norm on the actual cutoff domain.

The overlap, background-support, tunable-density and periodic
arguments below have passed separate internal mathematical audits.


% No preamble. Unique prefix oi-. Standard theorem environments and amsmath.
% Abstract overlapping-support Gaussian cutoff theorem; actual RG matching
% and strict background support are explicitly separate.

\section{Overlapping interactions in a Gaussian cutoff integral}
\label{oi-section}

The disjoint-cell cutoff expansion extends to overlapping interaction
supports, including supports of unbounded size, under an explicit
weighted multiplier norm. The essential step is to retain the original
cell endpoints when differentiating covariance edges between larger
interaction molecules. This controls the degree factorials without
enlarging the fixed Gaussian cell dimension. The statement below is
abstract; it does not verify the fluctuation coordinates or interaction
bounds of the Yang--Mills construction.

\subsection{Assumptions and quantitative statement}

Let $\Lambda$ be a finite metric set of cells, each carrying
$x_i\in\mathbb R^D$, with $D$ fixed. Let $A$ vary in an open complex
background domain. The complex symmetric covariance $C(A)$ is
holomorphic and obeys, uniformly in volume and background,
\begin{equation}\label{oi-covariance}
 \operatorname{Herm}C\ge\lambda I,\qquad
 \|C\|\le\Lambda_0,\qquad 0<\lambda\le\Lambda_0.
\end{equation}
For $i\ne j$ assume
\begin{equation}\label{oi-decay}
 \|C_{ij}\|\le J_{ij}=J_{ji},\qquad J_{ii}=0,\qquad
 A_\kappa=\sup_i\sum_j e^{\kappa d(i,j)}J_{ij}<\infty.
\end{equation}
Write $\ell(S)$ for the minimum metric spanning-tree length of a
nonempty set $S$, with $\ell(\{i\})=0$. Gaussian expectations use
\[
 p_C(x)=(2\pi)^{-D|\Lambda|/2}(\det C)^{-1/2}
                     \exp(-x^TC^{-1}x/2),
\]
with the analytic determinant branch on the accretive cone.

Fix the real cutoff
\[
 \chi_i=\mathbf1_{\{\|x_i\|_2\le R\}},\qquad
 \chi_Y=\prod_{i\in Y}\chi_i.
\]
Consider a finite labelled family of nonempty interaction supports
$Y\subset\Lambda$. Distinct labels may share a support.
The functions $V_Y(A,x_Y)$ are measurable in the real variables and
holomorphic in $A$ where $\chi_Y=1$, and
\begin{equation}\label{oi-multiplier}
 v_Y=\sup_{A,\;\chi_Y=1}|e^{V_Y(A,x_Y)}-1|<\infty.
\end{equation}
No values outside that cutoff region or real-field derivatives are
required. Define $w_Y=\chi_Y(e^{V_Y}-1)$, extended by zero.
Set
\begin{align}
 \delta&=\lambda/2,&
 V_*&=8\Lambda_0^2/\lambda,&
 K_G&=(2V_*/\lambda)^{D/2},\notag\\
 p_R&=2^{D/2}e^{-R^2/(4V_*)},&
 \tau&=K_Gp_R,& c_G&=e^2/\delta,                 \label{oi-constants}\\
 \epsilon&=\sup_i\sum_{Y\ni i}
      v_Y K_G^{|Y|}e^{4|Y|+\kappa\ell(Y)},&
 \eta&=e^3\tau+\epsilon,&
 \eta_*&=\frac1{4\max(1,c_GA_\kappa)}.           \label{oi-smallness}
\end{align}
The interaction sum counts labels. These sufficient constants are
not asserted to be optimal.

\begin{theorem}[Overlapping Gaussian cutoff interactions]
\label{oi-theorem}
If $\eta\le\eta_*$, then
\[
 Z_V(A)=\mathbb E_{C(A)}
   \left[\prod_{i\in\Lambda}\chi_i
              \exp\left(\sum_YV_Y\right)\right]
\]
is nonzero and has a holomorphic logarithm defined by a convergent
expansion
\begin{equation}\label{oi-log}
 \log Z_V(A)=\sum_{\varnothing\ne X\subset\Lambda}U_V(X;A),\qquad
 \sup_i\sum_{X\ni i}e^{\kappa\ell(X)+|X|}
       \|U_V(X;\cdot)\|_\infty\le2\eta .
\end{equation}
The constants do not depend on the number or sizes of supports.
If
\begin{equation}\label{oi-quarter}
 e^3\tau\le\eta_*/4,\qquad \epsilon\le\eta_*/4,
\end{equation}
the canonical expansion also satisfies
\begin{equation}\label{oi-difference}
 \sup_i\sum_{X\ni i}e^{\kappa\ell(X)+|X|}
      \|U_V(X;\cdot)-U_0(X;\cdot)\|_\infty
 \le8\epsilon ,
\end{equation}
where $U_0$ is the pure Gaussian cutoff forest/Mayer expansion.
The support label $X$ concerns fluctuation cells; strict locality
in the background is a separate assertion.
\end{theorem}

\subsection{First expansion: disjoint interaction molecules}

Put $q_i=-\mathbf1_{\mathrm{bad},i}$. The identity
\begin{equation}\label{oi-product}
 \prod_i\chi_i e^{\sum_YV_Y}
       =\prod_i(1+q_i)\prod_Y(1+w_Y)
\end{equation}
is exact, and $q_iw_Y=0$ whenever $i\in Y$.
For a nonempty set $S$ let
\begin{align}
 a_S(x_S)&=
 \sum_{\substack{\mathcal F\ne\varnothing,\;\cup_{Y\in\mathcal F}Y=S\\
                  \mathcal F\ {\rm overlap\text{-}connected}}}
                         \prod_{Y\in\mathcal F}w_Y(x_Y),\notag\\
 M_S&=
 \sum_{\substack{\mathcal F\ne\varnothing,\;\cup_{Y\in\mathcal F}Y=S\\
                  \mathcal F\ {\rm overlap\text{-}connected}}}
                         \prod_{Y\in\mathcal F}v_Y .       \label{oi-molecules}
\end{align}
Each label is used at most once in these sums. Connectivity refers
to the intersection graph of the chosen supports. Then
$|a_S|\le M_S$. Define $f_S=a_S$ for nonsingletons and
$f_{\{i\}}=a_{\{i\}}+q_i$. Grouping overlap components in
\eqref{oi-product} gives exactly
\begin{equation}\label{oi-disjoint}
 \prod_i\chi_i e^{\sum_YV_Y}
   =\sum_{\mathcal P\ {\rm pairwise\ disjoint}}
                               \prod_{S\in\mathcal P}f_S(x_S).
\end{equation}
The grouping is unique: interaction components and bad singleton
marks have disjoint supports. Introduce
\begin{equation}\label{oi-b}
 b(S)=K_G^{|S|}M_S+\tau\mathbf1_{\{|S|=1\}}.
\end{equation}

\begin{lemma}[Overlap-animal estimate]\label{oi-animal}
If $\epsilon\le1$, then
\begin{equation}\label{oi-primitive-norm}
 \sup_i\sum_{S\ni i}b(S)e^{\kappa\ell(S)+3|S|}
       \le e^3\tau+\epsilon=\eta .
\end{equation}
\end{lemma}
\begin{proof}
Set $u_Y=v_YK_G^{|Y|}e^{3|Y|+\kappa\ell(Y)}$.
For an overlap-connected family,
$\ell(\cup Y)\le\sum_Y\ell(Y)$ and
$|\cup Y|\le\sum_Y|Y|$: join spanning trees of the supports,
which share vertices along their intersection graph.
Since $K_G\ge1$, its contribution to the interaction part of
\eqref{oi-primitive-norm} is at most $\prod_Yu_Y$.

Majorize a connected family anchored at $i$ by rooted spanning
trees of its intersection graph, choosing a root support containing
$i$. Dropping distinctness of labels only increases the sum.
The nonnegative rooted-tree majorant is the monotone iterative
solution of
\[
 T_Y=u_Y\exp\left(\sum_{Z:\,Z\cap Y\ne\varnothing}T_Z\right).
\]
If $T_Z\le u_Ze^{|Z|}$, then
\[
 \sum_{Z:\,Z\cap Y\ne\varnothing}T_Z
 \le\sum_{j\in Y}\sum_{Z\ni j}u_Ze^{|Z|}
 \le\epsilon |Y|\le |Y|.
 \]
Thus the bound persists through iteration, and
$\sum_{Y\ni i}T_Y\le\epsilon$. Every original connected family
is included at least once. The remaining bad singleton term is
$e^3\tau$.
\end{proof}

\subsection{Covariance forests with original cell endpoints}

Fix disjoint primitive supports $S_1,\ldots,S_n$. Introduce a
covariance parameter $s_{ab}$ for each pair of molecules, multiplying
all blocks between $S_a$ and $S_b$ and leaving internal blocks
unchanged. Partition masks in these parameters are direct sums of
actual principal submatrices of $C$. Their convex hull retains
\eqref{oi-covariance}, so the BKAR forest formula applies there.
At a forest mask both the covariance and the product of primitive
functions factorize over the forest components.

For a molecule-tree $T$, expand each covariance derivative into
its original endpoint cells:
\begin{equation}\label{oi-cell-derivative}
 \partial_{s_{ab}}
      =\sum_{i\in S_a,\;j\in S_b}\partial_{z_{ij}},
\end{equation}
where $z_{ij}$ multiplies just $C_{ij}$ and $C_{ji}$.
Choose one endpoint pair for each tree edge and let $d_i$ count
the resulting degree of cell $i$. No chosen endpoint pair repeats.
At the real forest mask, allow complex increments with radii
\begin{equation}\label{oi-radii}
 r_{ij}=\frac{\delta}{\max(d_i,d_j)J_{ij}}.
\end{equation}
Zero covariance edges contribute zero and can be omitted.
Every block row and column of the perturbation has sum at most
$\delta$, hence its operator norm is at most $\delta$. On the
whole polydisc,
\[
 \operatorname{Herm}M\ge\lambda I/2,\quad
 \|M\|\le2\Lambda_0,\quad
 \operatorname{Herm}M^{-1}\ge V_*^{-1}I,\quad
 |\det M|\ge(\lambda/2)^{D|\cup S_a|}.
\]
Thus $|p_M|$ is bounded by $K_G^{|\cup S_a|}$ times the product
real Gaussian density with cell covariance $V_*I$. For singletons
use $|f_i|\le M_i+\mathbf1_{\mathrm{bad},i}$ and its independent
Gaussian tail bound $p_R$. For other supports use $|f_S|\le M_S$.
This proves
\[
 \left|\mathbb E_M\prod_a f_{S_a}\right|
       \le\prod_a b(S_a).
\]
Cauchy's estimate, once in every chosen covariance edge, yields
\begin{equation}\label{oi-tree-bound}
 \left|\partial_{T,\ {\rm endpoints}}
          \mathbb E_{C(s)}\prod_a f_{S_a}\right|
 \le c_G^{n-1}\prod_a b(S_a)
                   \prod_i d_i!\prod_{ij\in T}J_{ij}.
\end{equation}
Indeed,
$\prod_{ij}\max(d_i,d_j)\le\prod_i d_i^{d_i}
 \le e^{2(n-1)}\prod_i d_i!$.
The factorials are attached to original cells, including degree-zero
cells with $0!=1$, and not to the sizes or degrees of whole molecules.
Neither indicators nor interaction multipliers are differentiated.

Forest components now regroup \eqref{oi-disjoint} into the exact
hard-core gas
\begin{equation}\label{oi-effective-gas}
 Z_V=\sum_{\mathcal Q\ {\rm pairwise\ disjoint}}
                                 \prod_{X\in\mathcal Q}\zeta(X).
\end{equation}
Here $\zeta(X)$ sums partitions of $X$ into disjoint primitive
supports and molecule trees with their BKAR integrals. The
one-primitive case is its un-interpolated expectation.

\subsection{Endpoint factorials and the rooted-tree sum}

\begin{lemma}\label{oi-tree-lemma}
Under \eqref{oi-primitive-norm} and $\eta\le\eta_*$,
\begin{equation}\label{oi-zeta-bound}
 \sup_i\sum_{X\ni i}|\zeta(X)|
                  e^{\kappa\ell(X)+2|X|}\le\eta .
\end{equation}
\end{lemma}
\begin{proof}
Put $\widehat b(S)=b(S)e^{\kappa\ell(S)+2|S|}$ and
$\widehat J_{ij}=e^{\kappa d(i,j)}J_{ij}$.
The union weight is dominated by these support weights and endpoint
edge weights: spanning trees inside the supports together with
the chosen covariance edges connect their union.

Root a molecule tree and retain the incoming cell at each nonroot
vertex. Let $F_i$ majorize planted trees whose root molecule
contains $i$ and carries one extra incoming half-edge at $i$.
This half-edge has no covariance weight; it increments the cell
degree. Allowing repeated or overlapping primitive supports
enlarges the positive majorant. With
$t_k=c_G\sum_j\widehat J_{kj}F_j$, its recursion is
\begin{equation}\label{oi-recursion}
 F_i=\sum_{S\ni i}\widehat b(S)
              (1-t_i)^{-1}\prod_{k\in S}(1-t_k)^{-1}.
\end{equation}
To verify its normalization, suppose $n_k$ children depart from
cell $k$. The unordered-child factor, after summing endpoint
assignments, is $1/\prod_k n_k!$. The cell-degree factorial in
\eqref{oi-tree-bound} is $\prod_k n_k!$, except that the incoming
cell has $(n_i+1)!$. The residual factors are $1$ at every
ordinary cell and $n_i+1$ at the incoming cell. Their nonnegative
geometric series give exactly \eqref{oi-recursion}.
This can equivalently be obtained from ordered children weighted
by the inverse factorial of their total number. It is the rooted
labelled-tree exponential series; repeated labels are used only
for the upper bound.

Define the recursion by finite-height truncation. If $F_j\le2\eta$,
then $t_k\le2c_GA_\kappa\eta\le1/2$, and
\[
 F_i\le2\sum_{S\ni i}\widehat b(S)2^{|S|}
      \le2\sum_{S\ni i}\widehat b(S)e^{|S|}
      \le2\eta.
\]
This gives a convergent planted majorant. At an unplanted root
the extra factor $(1-t_i)^{-1}$ is absent, so its sum over root
supports containing $i$ is at most
$\sum_{S\ni i}\widehat b(S)2^{|S|}\le\eta$.
In an original disjoint primitive family whose union contains
$i$, precisely one support contains $i$; root the tree there.
This proves \eqref{oi-zeta-bound}.

In particular a molecule of $m$ cells with equal endpoint parameter
$t$ has unplanted generating function $(1-t)^{-m}$. Its degree-$d$
coefficient is $(m)_d/d!$, a rising factorial divided by $d!$.
Replacing this by the crude block bound $m^d d!$ would lose the
support-exponential control.
\end{proof}

\subsection{Logarithm and analytic comparison}

For $z'(X)=|\zeta(X)|e^{\kappa\ell(X)+|X|}$,
\eqref{oi-zeta-bound} gives the Kotecky--Preiss condition
with size function $|X|$:
\[
 \sum_{Y:\,Y\cap X\ne\varnothing}z'(Y)e^{|Y|}
       \le |X|\eta\le |X|.
\]
The absolutely convergent Mayer logarithm, grouped by the union
of each connected overlap cluster, defines $U_V$. The union
weight is at most the product of its polymer weights. Anchor a
cluster at $i$ by distinguishing a polymer containing $i$ and
apply the pinned cluster estimate. It gives
\begin{equation}\label{oi-pinned}
 \sup_i\sum_{X\ni i}e^{\kappa\ell(X)+|X|}|U_V(X)|
 \le\sup_i\sum_{Y\ni i}
       |\zeta(Y)|e^{\kappa\ell(Y)+2|Y|}\le\eta.
\end{equation}
Uniform background suprema can be taken at every preceding step.
Locally uniform convergence proves holomorphy and identifies
$Z_V$ with the exponential of this logarithm, proving
nonvanishing and \eqref{oi-log}.

For \eqref{oi-difference}, replace every $w_Y$ by $t w_Y$ while
retaining the bad singleton marks. This interpolates bounded
multipliers. Its interaction norm is $|t|\epsilon$.
If $\epsilon>0$, set $\rho=\eta_*/(2\epsilon)\ge2$.
Under \eqref{oi-quarter}, all estimates hold for $|t|\le\rho$
because $\eta(t)\le3\eta_*/4$. The canonical activity vector is
holomorphic on this disk in the rooted activity norm and bounded
by $2\eta_*$. Banach-valued Cauchy therefore gives
\[
 \|U(1)-U(0)\|\le\frac{2\eta_*}{\rho-1}\le8\epsilon .
\]
The case $\epsilon=0$ is immediate. At $t=0$ only bad singleton
primitives remain, so $U(0)$ is the original pure Gaussian cutoff
forest/Mayer expansion. This completes the proof.

\subsection{Scope of the abstract result}

The same conclusion holds for countably many labelled interactions
under \eqref{oi-smallness}, defining the integrand by the uniformly
convergent product of their bounded multipliers on the cutoff,
or equivalently by its finite-family limit. No convergence of the
chosen logarithms $V_Y$ themselves is asserted. If
$u_Y=\sup_{\rm cutoff}|V_Y|\le1$, then
$v_Y\le2u_Y$ is a sufficient conversion to the required norm.
Finite range is optional; the weighted summability is the actual
hypothesis. The factor $K_G^{|Y|}$ budgets the Gaussian density
cost but does not allow the original cell dimension $D$ to grow
with the regulator.

The argument establishes locality of support labels in fluctuation
space. A covariance depending on a remote background variable may
give every activity the same remote dependence, even if the
covariance is diagonal. Strict background support consequently
requires the separate local-precision/support-algebra hypotheses.
Actual Yang--Mills use also requires the stated independent cell
coordinates, uniform complex spectrum, fixed cutoff domain and
weighted interaction bounds. None is inferred here from a global
interaction stability estimate or from pointwise inverse bounds.

The imported expansion results are the BKAR forest formula on the
convex hull of partition masks (Abdesselam--Rivasseau,
arXiv:hep-th/9409094) and the Kotecky--Preiss hard-core polymer
theorem, \emph{Communications in Mathematical Physics} \textbf{103}
(1986), 491--498. The molecule grouping, cell-endpoint Cauchy
estimate and endpoint-factorial tree sum above supply the new
overlapping-support step.

% No preamble. Unique prefix ob-. Follows the oi- overlap theorem.
% Strict background support under the explicit local-precision hypotheses.

\section{Strict background support for overlapping interactions}
\label{ob-section}

The preceding expansion can also be made strictly local in its
background variables. This requires a background-independent baseline
and a small local precision decomposition, in addition to local
background dependence of the interactions. The next statement checks
these additional inputs explicitly.

\subsection{Baseline, local precision and constants}

Let $\Lambda$ be a finite subset of $\mathbb Z^d$ with the ambient
lattice graph metric and fixed cell dimension $D$. The background
domain is a product of complex coordinate domains. Write
\begin{equation}\label{ob-precision}
 P(A)=P_0+\sum_Y B_Y(A|_Y),\qquad C_0=P_0^{-1},
\end{equation}
where $C_0$ is complex symmetric and independent of every background variable, each
$B_Y$ is complex symmetric, has rows and columns supported in $Y$,
and reads only $A|_Y$. Assume
\begin{equation}\label{ob-spectrum}
 \operatorname{Herm}C_0\ge\lambda I,\qquad
 \|C_0\|\le\Lambda_0,\qquad 0<\lambda\le\Lambda_0,
\end{equation}
and, for fixed $0<\kappa<m$,
\begin{align}
 K&=\sup_b\sum_c e^{(\kappa+2m)d(b,c)}
                         \|(C_0)_{bc}\|<\infty,\notag\\
 \nu&=\sup_b\sum_{Y\ni b}\sum_{c\in Y}
       e^{(\kappa+2m)\ell(Y)}
                         \|(B_Y)_{bc}\|_\infty<\infty. \label{ob-norms}
\end{align}
Symmetry gives the same column bounds. All constants are uniform
in volume and background. Define
\begin{align}
 \delta&=\lambda/4,&
 V_*&=8\Lambda_0^2/\lambda,&
 J_{bc}&=K e^{-m d(b,c)}\quad(b\ne c),\notag\\
 A_\kappa&=K\sum_{z\ne0}e^{-(m-\kappa)|z|},&
 \bar K&=K+\delta,& a&=\lambda/2,\notag\\
 s&=\min\{m/2,\;aem/(4\bar K)\},&
 L&=\frac4\lambda\sum_{z\in\mathbb Z^d}e^{-s|z|}. \label{ob-baseline-constants}
\end{align}
Set $J_{bb}=0$. Require $K\nu<1$ and put
\begin{equation}\label{ob-perturbation-constants}
 v=\frac{K^2\nu}{1-K\nu},\qquad
 v_1=v(1+\delta/K),\qquad
 u=Lv_1,\qquad w=\frac{L^2v_1}{1-u}.
\end{equation}
The sufficient precision smallness conditions are
\begin{equation}\label{ob-precision-smallness}
 v_1\le\lambda/4,\qquad u\le1/2,\qquad w\le1/(2V_*).
\end{equation}
They are chosen before increasing the cutoff radius.
The covariance is $C(A)=P(A)^{-1}$, given by its convergent
Neumann series.

Retain the fixed cutoff and interaction multiplier hypotheses
of \eqref{oi-multiplier}, now requiring each $V_Y$ to read
background only in $Y$. Define
\begin{align}
 \widehat K_G&=(8V_*/\lambda)^{D/2},&
 \widehat p_R&=2^{D/2}e^{-R^2/(8V_*)},&
 \widehat\tau&=\widehat K_G\widehat p_R,\notag\\
 \widehat\epsilon&=
  \sup_i\sum_{Y\ni i}
   v_Y\widehat K_G^{|Y|}e^{4|Y|+\kappa\ell(Y)},&
 \widehat\eta&=e^3\widehat\tau+\widehat\epsilon,&
 \widehat\eta_*&=\frac1{4\max(1,e^2A_\kappa/\delta)}.
                                                        \label{ob-smallness}
\end{align}

\begin{theorem}[Strictly local overlapping-interaction activities]
\label{ob-theorem}
Under \eqref{ob-precision}--\eqref{ob-precision-smallness} and
$\widehat\eta\le\widehat\eta_*$, the logarithm has a convergent
representation by connected lattice supports:
\begin{equation}\label{ob-local-log}
 \log Z_V(A)=\sum_H E_{\Lambda,H}(A|_H),\qquad
 \sup_b\sum_{H\ni b}e^{a_*|H|}
        \|E_{\Lambda,H}\|_\infty
 \le2S_{\kappa/2}\widehat\eta,
\end{equation}
where
\[
 a_*=\min(1/2,\kappa/4),\qquad
 S_{\kappa/2}=\sum_{z\in\mathbb Z^d}e^{-(\kappa/2)|z|}.
\]
Each activity is holomorphic and reads only the indicated
background coordinates. Supports may lie in the lattice hull of
$\Lambda$. If
\begin{equation}\label{ob-quarter}
 e^3\widehat\tau\le\widehat\eta_*/4,\qquad
 \widehat\epsilon\le\widehat\eta_*/4,
\end{equation}
the same canonical choice obeys
\begin{equation}\label{ob-difference}
 \sup_b\sum_{H\ni b}e^{a_*|H|}
  \|E_{\Lambda,H,V}-E_{\Lambda,H,0}\|_\infty
 \le8S_{\kappa/2}\widehat\epsilon.
\end{equation}
\end{theorem}

\subsection{The support algebra and the actual marginal}

For a nonempty moment support $M$, let $\ell_M(X)$ be the
minimum total length of a forest on $M\cup X$ whose every
component meets $M$. The roots may remain in different components.
For formal local families $F=(F_X)_X$, define
\begin{equation}\label{ob-algebra}
 \|F\|_{\mathcal B_M}
   =\sum_Xe^{\kappa\ell_M(X)}\|F_X\|_\infty,\qquad
 (FG)_Z=\sum_{X\cup X'=Z}F_XG_{X'}.
\end{equation}
Every $F_X$ is holomorphic in $A|_X$. The product-domain
assumption gives its uniform local-function norm.
The inequality
$\ell_M(X\cup X')\le\ell_M(X)+\ell_M(X')$
makes this a Banach algebra. Enlarging the root set is
contractive. Evaluation, obtained by summing the local
coefficients, is a bounded algebra homomorphism.

Expand $C=C_0-C_0BC_0+C_0BC_0BC_0-\cdots$ in the individual
precision labels $B_Y$. Retain these labels before grouping by
the union of their supports. For endpoints $b,c\in M$, every
term connects all its background labels to its endpoints by
covariance links and trees inside the supports $Y$. If their
total length is $L_{\rm path}$, both
$\ell_M(\cup Y)\le L_{\rm path}$ and
$d(b,c)\le L_{\rm path}$. Successive weighted matrix sums give
$K^{j+1}\nu^j$ at order $j$. Therefore, for
$\Delta C=C-C_0$,
\begin{equation}\label{ob-covariance-algebra}
 \max\left\{
  \sup_{b\in M}\sum_{c\in M}e^{2md(b,c)}
                         \|(\Delta C)_{bc}\|_{\mathcal B_M},
  \sup_{c\in M}\sum_{b\in M}e^{2md(b,c)}
                         \|(\Delta C)_{bc}\|_{\mathcal B_M}
 \right\}\le v.
\end{equation}
This bounds the whole covariance marginal, including its
background dependence, rather than only visible forest edges.

\subsection{Molecule masks and cell endpoint polydiscs}

At a molecule BKAR point, its induced cell mask $\sigma$
is one inside every primitive molecule and has the BKAR values
between molecules. Each molecule partition mask is exactly
a cell partition mask whose blocks are unions of molecules.
Thus $\sigma$ is in the convex hull of cell partition masks
and $0\le\sigma_{bc}\le1$.

Choose one cell endpoint pair for each molecule-tree edge,
with cell degrees $d_b$, and use
$|z_{bc}|<\delta/[\max(d_b,d_c)J_{bc}]$.
The endpoint graph need not span or connect every cell.
Let $M_0$ be the resulting masked and perturbed baseline.
The row-sum argument gives
\begin{equation}\label{ob-m0}
 \operatorname{Herm}M_0\ge\lambda I/2,\qquad
 \|M_0\|\le2\Lambda_0,\qquad
 \|M_0\|_{\mathrm{Schur},m}\le K+\delta.
\end{equation}
For the last bound, each selected edge satisfies
\[
 e^{md(b,c)}|z_{bc}|\|(C_0)_{bc}\|
 \le\frac{\delta}{\max(d_b,d_c)}
                     e^{-\kappa d(b,c)}.
\]
Its row sum is at most $\delta$, using only the number of
incident selected edges. No spanning-tree property on the
cell set has been used.

Conjugation by $D_{bb}=e^{s d(b,c)}$ changes $M_0$ in
operator norm by at most
$2s(K+\delta)/(em)\le a/2$. Since its Hermitian part is at
least $a$, the inverse entries are bounded by
$(2/a)e^{-sd(b,c)}$, and hence
\begin{equation}\label{ob-inverse-schur}
 \|M_0^{-1}\|_{\mathrm{Schur}}\le L.
\end{equation}

Apply the same mask and increments to $\Delta C$, obtaining
an algebra-valued matrix $E$. The mask part has norm at most
$v$. By \eqref{ob-covariance-algebra}, the selected-edge part
has row and column sums at most $\delta v/K$:
\[
 |z_{bc}|v e^{-2md(b,c)}
   \le\frac{\delta v}{K\max(d_b,d_c)}e^{-md(b,c)}.
\]
Consequently
\begin{equation}\label{ob-e}
 \|E\|_{\mathrm{Schur}(\mathcal B_M)}\le v_1.
\end{equation}
These are precisely the estimates needed for the density
expansion, and they remain valid for the molecule masks and
the original cell endpoint polydiscs.

\subsection{The entire density, including determinant normalization}

Put $F=M_0^{-1}E$. In the matrix support algebra,
\[
 (M_0+E)^{-1}-M_0^{-1}
       =\sum_{j\ge1}(-F)^jM_0^{-1},
\]
whose Schur norm is at most $w$, while
\[
 \|\operatorname{Tr}\log(I+F)\|_{\mathcal B_M}
       \le D|M|[-\log(1-u)].
\]
The exact normalized density identity is
\begin{equation}\label{ob-density-identity}
 p_{M_0+E}(x)=p_{M_0}(x)
 \exp\left[
 -\tfrac12\operatorname{Tr}\log(I+F)
 -\tfrac12x^T\big((M_0+E)^{-1}-M_0^{-1}\big)x
 \right].
\end{equation}
Taking its algebra norm, the determinant factor costs at most
$2^{D|M|/2}$. The quadratic factor costs at most
$e^{w\|x\|^2/2}$, which by
\eqref{ob-precision-smallness} uses at most half of the
baseline Gaussian decay. It follows that
\begin{equation}\label{ob-density-bound}
 \|p_{M_0+E}(x)\|_{\mathcal B_M}
       \le\widehat K_G^{|M|}p_{2V_*I}(x_M).
\end{equation}
Here $p_{2V_*I}$ is the independent product Gaussian density
with that covariance in every original cell. All inverse,
logarithm, exponential and support series are absolutely
convergent under the displayed majorants.

For an interaction supported on $Y\subset M$,
$\ell_M(Y)=0$. Thus $w_Y$ is an algebra element with
norm at most $v_Y$. A primitive interaction molecule on $S$
has norm at most $M_S$ from \eqref{oi-molecules}; a bad
singleton has empty background support.
Integrating \eqref{ob-density-bound} against their product
gives on the entire endpoint polydisc
\begin{equation}\label{ob-moment-bound}
 \left\|\mathbb E_{M_0+E}\prod_a f_{S_a}\right\|_{\mathcal B_M}
 \le\prod_a\left(
   \widehat K_G^{|S_a|}M_{S_a}
       +\widehat\tau\mathbf1_{\{|S_a|=1\}}\right).
\end{equation}

\subsection{Factorization with overlapping background labels}

Retain independent formal scalar tags on the precision labels,
and on the interaction labels if desired, until absolute
convergence is proved. At a forest mask the complete marginal
Gaussian matrix is block diagonal over the forest components.
Both its determinant and inverse therefore factorize in
\eqref{ob-density-identity}, giving an identity of formal
power series. This remains true when different components
read overlapping background coordinates. Grouping coefficients
by the union of their background labels respects multiplication.
The density for a component is its actual marginal density,
with its own root set; injection into the algebra with all
moment roots is contractive.

Thus forest factorization holds in the support representation
itself, not merely after evaluating $A$.
Banach-valued Cauchy applied to \eqref{ob-moment-bound} yields
the same cell-factorial bound \eqref{oi-tree-bound}.
The overlap-animal and planted-tree estimates give
\begin{equation}\label{ob-zeta}
 \sup_i\sum_{M\ni i,X}
 e^{\kappa\ell(M)+\kappa\ell_M(X)+2|M|}
            \|\zeta(M;X)\|_\infty\le\widehat\eta.
\end{equation}
Hard-core compatibility is still disjointness of the marked
supports $M$, and does not require disjoint background supports.

For a connected Mayer cluster, the root and union inequalities
give
\[
 \ell(\cup M_a)\le\sum_a\ell(M_a),\qquad
 \ell_{\cup M_a}(\cup X_a)
                  \le\sum_a\ell_{M_a}(X_a).
\]
The same pinned cluster estimate therefore yields
\begin{equation}\label{ob-marked-log}
 \log Z_V=\sum_{M\ne\varnothing,X}U(M;X)(A|_X),\qquad
 \sup_i\sum_{M\ni i,X}
 e^{\kappa[\ell(M)+\ell_M(X)]+|M|}
                 \|U(M;X)\|_\infty\le2\widehat\eta.
\end{equation}

\subsection{Re-rooting at the final connected background support}

Choose deterministic minimizing trees for $M$, minimizing
forests attaching $X$ to $M$, and shortest lattice paths for
their edges. Their union $H(M,X)$ is connected and contains
$M\cup X$. With $L_*=\ell(M)+\ell_M(X)$,
\[
 |H(M,X)|\le |M|+L_*,\qquad
 d(a,b)\le L_*\quad(a\in M,\ b\in H(M,X)).
\]
Consequently
\begin{equation}\label{ob-hull}
 e^{a_*|H|}
 \le e^{|M|+\kappa L_*}e^{-(\kappa/2)d(a,b)}.
\end{equation}
Group $U(M;X)$ by this fixed hull to define $E_{\Lambda,H}$.
To root a sum at $b\in H$, overcount by all $a\in M$ and
use \eqref{ob-marked-log}; the lattice convolution costs
only $S_{\kappa/2}$. This proves \eqref{ob-local-log}.

Under \eqref{ob-quarter}, the multiplier interpolation
$w_Y\mapsto tw_Y$ is holomorphic in the same support algebra.
The Banach-valued Cauchy estimate of the preceding section
gives $8\widehat\epsilon$ in the marked rooted norm.
The fixed hull map has norm at most $S_{\kappa/2}$ by
\eqref{ob-hull}, proving \eqref{ob-difference}.

\subsection{What the locality theorem does and does not identify}

The activities may depend on the whole background-independent
finite-volume baseline $C_{0,\Lambda}$. Their variable background
dependence reads only $A|_H$, but volume consistency is not
claimed. A baseline chosen separately at each arbitrary
background would localize only perturbations from that reference.
The product background domain and the small local precision
norm are substantive hypotheses, not consequences of a weak
global gauge-field norm or of covariance decay alone.

A deterministic path choice may preserve translations while
breaking discrete rotations; arbitrary decompositions do not
make each activity gauge invariant. An application needing
these symmetries must supply an equivariant construction or
an admissible averaging argument.
The periodic extension later in this notebook verifies uniform
volume counting and the required path/hull construction with the
same length and cardinality bounds, including short tori.
No actual regulator coordinates, source coefficients, moving
gauge chart or admitted history family has been identified by
the abstract hypotheses alone.

% No preamble. Unique prefix otg-. Follows oi- and ob-.
% Real-baseline refinement with tunable density and joint-tail constants.

\section{A tunable density cost near a real Gaussian baseline}
\label{otg-section}

The fixed density factors in the preceding theorems can cost
$\exp(cD)$ per interaction cell. That cost need not fit an
application's support-decay exponent when a cell contains many
coordinates. A real baseline allows a tunable factor
$\exp(O(\theta D))$. The dominating Gaussian is correlated;
a joint-tail estimate replaces independence.

\begin{lemma}[Relative Gaussian density]\label{otg-density-lemma}
Let $C_r$ be a real symmetric positive covariance in $N$
dimensions, and let
\[
 M=C_r^{1/2}(I+F)C_r^{1/2},\qquad F^T=F,\qquad
                       \|F\|\le\theta<1.
\]
Then
\begin{equation}\label{otg-density}
 |p_M(x)|
 \le\left(\frac{1+\theta}{1-\theta}\right)^{N/2}
                                p_{(1+\theta)C_r}(x).
\end{equation}
\end{lemma}
\begin{proof}
The determinant modulus obeys
$|\det(I+F)|\ge(1-\theta)^N$. Moreover
\[
 (1+\theta)\operatorname{Herm}(I+F)-(I+F)^*(I+F)
 =\theta I-(1-\theta)\operatorname{Herm}F-F^*F\ge0,
\]
using $\operatorname{Herm}F\le\theta I$ and
$F^*F\le\theta^2I$.
Congruence by $(I+F)^{-1}$ gives
$\operatorname{Herm}(I+F)^{-1}\ge(1+\theta)^{-1}I$.
Insert these two bounds in the Gaussian density, including
the normalization of the real density on the right.
\end{proof}

\begin{lemma}[Correlated block tails]\label{otg-tail-lemma}
Let $G$ be a centered real Gaussian on $D$-dimensional cells, with
covariance at most $\Lambda_{\rm g} I$. For $0<\sigma\le1$ set
\begin{equation}\label{otg-tail-constant}
 p_{\sigma,R}(\Lambda_{\rm g})
  =(1+\sigma)^{D/2}
             \exp\left[-\frac{\sigma R^2}{2(1+\sigma)\Lambda_{\rm g}}\right].
\end{equation}
For every finite cell set $B$,
\begin{equation}\label{otg-joint-tail}
 \mathbb P\{\|G_i\|_2>R\text{ for all }i\in B\}
                   \le p_{\sigma,R}(\Lambda_{\rm g})^{|B|}.
\end{equation}
Consequently, for $a_i\ge0$,
\begin{equation}\label{otg-product-tail}
 \mathbb E\prod_{i\in I}
            (a_i+\mathbf1_{\{\|G_i\|_2>R\}})
                  \le\prod_{i\in I}(a_i+p_{\sigma,R}(\Lambda_{\rm g})).
\end{equation}
\end{lemma}
\begin{proof}
On the event in \eqref{otg-joint-tail},
$\sum_{i\in B}\|G_i\|^2>|B|R^2$.
Exponential Markov with
$t=\sigma/[2(1+\sigma)\Lambda_{\rm g}]$ gives
\[
 \mathbb P(\text{event})
 \le e^{-t|B|R^2}\det(I-2tC_B)^{-1/2}
 \le e^{-t|B|R^2}(1+\sigma)^{D|B|/2}.
\]
Expand the product into choices of bad singleton cells to
obtain \eqref{otg-product-tail}. No independence is assumed.
\end{proof}

\subsection{The marked theorem for a small complex covariance family}

Let $C_r$ be a fixed background-independent real covariance with
$\lambda I\le C_r\le\Lambda_0 I$. Suppose the actual holomorphic
complex symmetric covariance satisfies
\begin{equation}\label{otg-close}
 \sup_A\|C(A)-C_r\|\le\theta\lambda/2,\qquad
                                  0<\theta\le1/4.
\end{equation}
Retain the actual covariance majorant $J$ and $A_\kappa$
from \eqref{oi-decay}, the fixed cutoff, and the bounded local
interaction multipliers of the marked theorem.
For any fixed $0<\sigma\le1$ replace its constants by
\begin{align}
 \delta_\theta&=\theta\lambda/2,&
 K_\theta&=\left(\frac{1+\theta}{1-\theta}\right)^{D/2}
                     =e^{D\operatorname{artanh}\theta},\notag\\
 p_{\theta,\sigma,R}
       &=p_{\sigma,R}((1+\theta)\Lambda_0),&
 \tau_\theta&=K_\theta p_{\theta,\sigma,R},&
 c_\theta&=e^2/\delta_\theta,\notag\\
 \epsilon_\theta&=\sup_i\sum_{Y\ni i}
       v_YK_\theta^{|Y|}e^{4|Y|+\kappa\ell(Y)},&
 \eta_\theta&=e^3\tau_\theta+\epsilon_\theta,&
 \eta_{\theta,*}&=\frac1{4\max(1,c_\theta A_\kappa)}.
                                                        \label{otg-mark-constants}
\end{align}

\begin{theorem}[Tunable marked expansion]\label{otg-mark-theorem}
If $\eta_\theta\le\eta_{\theta,*}$, the conclusion
\eqref{oi-log} holds with $2\eta_\theta$.
Under the quarter budgets
$e^3\tau_\theta\le\eta_{\theta,*}/4$ and
$\epsilon_\theta\le\eta_{\theta,*}/4$,
the interaction-difference bound is $8\epsilon_\theta$.
\end{theorem}
\begin{proof}
A cell partition mask is a direct-sum pinching, and therefore
contracts the operator norm of $C(A)-C_r$. Convex combinations
of these masks have the same property. At a molecule BKAR
mask the covariance differs from its real masked baseline by
at most $\theta\lambda/2$. The degree-adapted cell endpoint
polydisc, with radius budget $\delta_\theta$ and the actual
majorant $J$, adds at most another $\theta\lambda/2$.
The total relative perturbation from the real masked baseline
has norm at most $\theta$.

Lemma~\ref{otg-density-lemma} applies. Its dominating real
Gaussian covariance is at most $(1+\theta)\Lambda_0I$.
Bound nonsingleton molecule functions by $M_S$, and expand
the singleton bounds $M_i+\mathbf1_{\mathrm{bad},i}$.
Lemma~\ref{otg-tail-lemma} gives exactly the product bound
used in the cell-endpoint Cauchy estimate, with
\[
 b_\theta(S)=K_\theta^{|S|}M_S+
                      \tau_\theta\mathbf1_{\{|S|=1\}}.
\]
The overlap-animal, planted-tree, Mayer and interpolation
arguments are unchanged.
\end{proof}

If $C(A)$ is identically the real baseline, one may use the
larger budget $\delta_\theta=\theta\lambda$, because no
background deviation consumes half of it.
Operator closeness suffices for this marked theorem but does
not imply strict background locality.

\subsection{The tunable strictly local version}

Assume the local precision and product-domain hypotheses of
Section~\ref{ob-section}, now with $C_0$ real symmetric positive.
Retain $K,\nu,v$, $J_{bc}=Ke^{-md(b,c)}$ and $A_\kappa$.
Fix $0<\theta\le1/4$ and replace the precision constants by
\begin{align}
 \delta&=\theta\lambda,& \bar K&=K+\delta,& a&=\lambda/2,
 \notag\\
 s&=\min\{m/2,\;aem/(4\bar K)\},&
 L&=\frac4\lambda\sum_z e^{-s|z|},\notag\\
 v_1&=v(1+\delta/K),&
 u&=Lv_1,& w&=\frac{L^2v_1}{1-u}.                \label{otg-precision-constants}
\end{align}
Require $K\nu<1$ and the strengthened smallness conditions
\begin{equation}\label{otg-precision-small}
 u\le\theta,\qquad
 w\le\frac{\theta}{(1+\theta)\Lambda_0}.
\end{equation}
The earlier condition $v_1\le\lambda/4$ follows because
$L\ge4/\lambda$ and $\theta\le1/4$.
Define
\begin{equation}\label{otg-strict-density-constants}
 \rho_\theta=\frac{1+\theta}{1-\theta},\qquad
 \mathcal K_\theta=
     \left(\frac{1+\theta}{(1-\theta)^3}\right)^{D/2},
 \qquad
 \mathcal p_{\theta,\sigma,R}
        =p_{\sigma,R}(\rho_\theta\Lambda_0).
\end{equation}
In \eqref{ob-smallness} replace
$\widehat K_G,\widehat p_R$ by
$\mathcal K_\theta,\mathcal p_{\theta,\sigma,R}$,
and use $c_G=e^2/\delta$ with the present $\delta$.
Denote the resulting quantities by
$\mathcal\tau_\theta,\mathcal\epsilon_\theta,
\mathcal\eta_\theta,\mathcal\eta_{\theta,*}$.

\begin{theorem}[Tunable strictly local expansion]\label{otg-strict-theorem}
Under \eqref{otg-precision-small} and
$\mathcal\eta_\theta\le\mathcal\eta_{\theta,*}$,
\eqref{ob-local-log} holds with bound
$2S_{\kappa/2}\mathcal\eta_\theta$.
Under the corresponding quarter budgets, its interaction
difference is bounded by
$8S_{\kappa/2}\mathcal\epsilon_\theta$.
All domain, volume and symmetry scope statements of
Section~\ref{ob-section} remain in force.
\end{theorem}
\begin{proof}
Let $C_{r,M}$ be the real masked baseline on a moment set $M$.
It is between $\lambda I$ and $\Lambda_0I$.
Add only its selected complex cell endpoint perturbations to
form $M_0$. Their norm is at most $\theta\lambda$, so
\[
 |p_{M_0}(x)|
          \le K_\theta^{|M|}p_{(1+\theta)C_{r,M}}(x).
\]
The support-algebra bounds on the source-dependent matrix
$E$ from Section~\ref{ob-section} are unchanged with the
present constants. In particular the determinant correction
has norm at most $(1-u)^{-D|M|/2}$ and the inverse correction
costs at most $\exp(w\|x\|^2/2)$.

Since $C_{r,M}\le\Lambda_0I$,
$\|x\|^2\le\Lambda_0x^TC_{r,M}^{-1}x$.
Condition \eqref{otg-precision-small} therefore gives
\[
 p_{(1+\theta)C_{r,M}}(x)e^{w\|x\|^2/2}
 \le\left(\frac{\rho_\theta}{1+\theta}\right)^{D|M|/2}
                                  p_{\rho_\theta C_{r,M}}(x).
\]
Combining this inequality with $u\le\theta$ and the
baseline density factor proves
\begin{equation}\label{otg-algebra-density}
 \|p_{M_0+E}(x)\|_{\mathcal B_M}
       \le\mathcal K_\theta^{|M|}
                                  p_{\rho_\theta C_{r,M}}(x).
\end{equation}
The dominating Gaussian is real and correlated, with
covariance at most $\rho_\theta\Lambda_0I$.
Expand the bad singleton choices and use
Lemma~\ref{otg-tail-lemma}. This gives
\eqref{ob-moment-bound} with the new constants.
All cell-endpoint, formal-factorization, Mayer, hull and
bounded-multiplier interpolation arguments then apply.
\end{proof}

\subsection{The dimension tradeoff}

For $\theta\le1/4$,
\[
 \log K_\theta\le\frac{16}{15}D\theta,\qquad
 \log\mathcal K_\theta
 =\frac D2[\log(1+\theta)-3\log(1-\theta)]
                   \le\frac52D\theta .
\]
Thus $\theta=1/(4D)$ bounds these unmarked per-cell factors
by $e^{4/15}$ and $e^{5/8}$, independently of $D$.
The covariance-edge constant grows like $\theta^{-1}$,
and the local precision neighborhood must meet
\eqref{otg-precision-small}. If the other baseline constants
stay fixed, this exchanges an exponential-in-$D$ interaction
support cost for a polynomial edge-budget loss and a smaller
precision neighborhood.

The choice $\sigma=1$ retains the faster Gaussian tail
$2^{D/2}e^{-R^2/(4\Lambda_{\rm g})}$. Alternatively,
$\sigma=1/(4D)$ bounds its prefactor by $e^{1/8}$,
but its exponent is only of order $R^2/(D\Lambda_{\rm g})$.
The cutoff must satisfy the resulting tail budget in either
case. The lemma does not show that the actual source
covariance, coordinate family, local precision or cutoff
admits these tuned neighborhoods. It removes a specific
artificial density cost while keeping those hypotheses
explicit.

% No preamble. Unique prefix op-. Follows the oi-, ob-, and otg- modules.

\section{Periodic volumes and short periods}
\label{op-section}

The preceding strict-support expansions also hold on finite rectangular
periodic lattices. This is a geometric extension of the abstract theorems;
their covariance, precision, interaction and background-domain assumptions
must still hold for the actual periodic data.

\begin{proposition}[Periodic strict-support expansion]
\label{op-periodic-theorem}
Let
\[
 \mathbb T_N=\prod_{\mu=1}^{d}\mathbb Z/N_\mu\mathbb Z,
 \qquad N_\mu\ge1,
\]
with its nearest-neighbor graph distance $d_N$. One may replace the finite
subset of $\mathbb Z^d$ in the strictly local overlap theorem and its
tunable real-baseline version by $\mathbb T_N$, or by any subset of it.
Use $d_N$ in all weighted matrix norms and in the metric tree and rooted
forest lengths. If the displayed hypotheses hold with constants independent
of $N$, both conclusions hold with exactly the same constants. The resulting
supports are connected subsets of $\mathbb T_N$, including when some
periods equal one or two.
\end{proposition}

\begin{proof}
For every base point $b\in\mathbb T_N$, choose one shortest integer
representative $r_b(c)\in\mathbb Z^d$ of each difference $c-b$.
The choices form an injection and satisfy
$|r_b(c)|_1=d_N(b,c)$. Consequently, for every nonnegative function
$f$ on the nonnegative integers,
\begin{equation}\label{op-counting}
 \sum_{c\in\mathbb T_N}f(d_N(b,c))
 \le\sum_{z\in\mathbb Z^d}f(|z|_1).
\end{equation}
For sums excluding $c=b$, the zero representative is excluded on the
right. In particular, all exponential lattice-sum bounds used for
$A_\kappa$, the inverse Schur constant, and $S_{\kappa/2}$ retain their
previous upper bounds.

The weighted matrix product estimates require only the triangle inequality.
The inverse Schur proof uses the weight
$D_{bb}=\exp(s d_N(b,c))$ and the inequality
\begin{equation}\label{op-weight-ratio}
 |d_N(b,c)-d_N(b',c)|\le d_N(b,b').
\end{equation}
Thus its conjugation estimate, operator inverse bound and entry decay are
unchanged; summing entries uses \eqref{op-counting}. Partition masks,
cell endpoint degrees, Gaussian density comparisons, the rooted support
algebra, and the overlap and Mayer combinatorics are independent of any
choice of integer representatives. They therefore apply directly on the
finite quotient.

It remains to check the conversion to connected background support. For
each nonempty marked set $M$ and background support $X$, let
$\ell_N(M)$ be its minimum metric tree length, and let
$\ell_{N,M}(X)$ be the minimum forest length on $M\cup X$ with every
component meeting $M$. Realize fixed minimizing trees and forests by
fixed shortest torus paths, and let $H(M,X)$ be their union. It is a
connected subset of $\mathbb T_N$, contains $M\cup X$, and, writing
$L_* =\ell_N(M)+\ell_{N,M}(X)$, satisfies
\begin{equation}\label{op-hull}
 |H(M,X)|\le |M|+L_*,\qquad
 d_N(a,b)\le L_*
 \quad(a\in M,\ b\in H(M,X)).
\end{equation}
The vertex bound follows by adding the selected paths one edge at a time;
identifications caused by wrapping can only reduce the number of vertices.
If the original cell set is a proper subset of the torus, the hull may
leave that set, while the coefficient continues to read only the original
background coordinates in $X$.

For $a_* =\min(1/2,\kappa/4)$, \eqref{op-hull} gives
\begin{equation}\label{op-reroot}
 e^{a_*|H(M,X)|}
 \le e^{|M|+\kappa L_*}
             e^{-(\kappa/2)d_N(a,b)}.
\end{equation}
Overcounting by the marked anchors $a\in M$ and applying
\eqref{op-counting} costs at most the same $S_{\kappa/2}$ as before.
This proves both the connected-support norm and the interaction-difference
bound with the original constants.
\end{proof}

The proof does not use a positive injectivity radius or distinct positive
and negative neighbors. Repeated interaction labels whose supports become
equal after a periodic identification are still counted as distinct labels,
as in the original multiplier norm. The theorem assumes the corresponding
actual periodic coefficient norms; it does not infer them merely by
relabeling an infinite-volume model. Deterministic shortest-path tie choices
need not preserve translations or rotations, and no such activity symmetry
is asserted here. Nor does this extension assert consistency between
different periodic volumes.

% No preamble. Source bridge for Notebook18. Unique prefix ib-.
\section{Using the published localization and integration theorem}
\label{ib-section}

There is an existing specialized integration route in the literature.
Ba\l aban's \emph{Renormalization group approach to lattice gauge
field theories II}, Commun. Math. Phys. \textbf{116} (1988),
Lemma 2, pp.~10--11, equations (1.41)--(1.43), localizes the
complete fluctuation exponent, including the difference of inherited
effective interactions. Lemma 3, p.~20, equation (2.38), and the
logarithmic estimate (2.41), p.~21, integrate these overlapping
terms. These results can be used in their published setting;
another abstract integration theorem is not a prerequisite for
using them.

\subsection{The actual coefficient and field domains}

The original domain is equation (1.34), p.~9:
\[
 \mathcal U^c_{k+1}
   (Y,(1+\beta)\alpha_0,(1+\beta)\alpha_1,\alpha_0)
       \;\times\;
 \{B:|B|<\epsilon_1/g\text{ on }Y\}.
\]
The first factor is an orbit domain from CMP109, pp.~262--263,
equations (1.10)--(1.16). Its restrictions include reconstructed
backgrounds on subregions. It cannot be replaced solely by an
unverified global supremum condition. The operator and minimizer
walk deformations in CMP116, pp.~3--6, are also inputs.

To distinguish two scales, let $M_{\rm loc}$ denote the fixed
localization cube size in that paper; it is different from the
refinement $M=L^k$ in our flat covariance theorems. The local
exponent is organized as
\[
 P^{(k)}+\text{inherited-action difference}=\sum_Y V_Y,
 \qquad
 V_Y=\tfrac12\langle B,Q_Y(U,J,B)B\rangle+V_Y''(U,J,B).
\]
The displayed source bounds are (1.43) for the quadratic kernel
and (1.36) for the remainder. Their localization weight is a
normalized minimum tree length meeting the $M_{\rm loc}$ cubes
of $Y$. It is stronger than a diameter bound and comes with
the source's specific geometric summations. The source exponent
contains the Haar, constraint Jacobian and gauge contributions;
it is not only the bare Wilson cubic remainder.

The original article was inspected through its extracted text.
Some displayed numerical coefficients are lost in that extraction;
the numbered source bounds are therefore located here rather than
reconstructed by guesswork. The matching deposited formulas are
Raw I:30590--30597 and Raw II:3190--3208. The latter attributes
(1.43) to CMP109, but that formula in this argument is CMP116
(1.43), as also recorded at Raw II:11178--11188.

\subsection{The logarithmic-radius transfer}

Fix all geometric and analytic parameters before choosing the
coupling ceiling. For $R(g)=c(\log g^{-2})^p$, $p>1/2$,
\[
 gR(g)<\epsilon_1/2,
 \qquad gR(g)^3\longrightarrow0,
 \qquad e^{-c_1R(g)^2}=o(g^q)
\]
for every fixed $q>0$. Thus cutoff points in the actual source
variable lie inside its fixed physical-field analytic radius.
For localized coefficients holomorphic on that radius and
vanishing at the indicated orders, radial Schwarz estimates give
\[
 \|Q_Y(gB)\|\le a_Y gR,
 \qquad |V_Y''(gB)|\le a_Y' gR.
\]
Fixed inverse powers of the analytic radius are included in
$a_Y,a_Y'$. On the cutoff this bounds $V_Y$ by $gR^3$ times
a local coefficient and a polynomial support factor. A reserved
part of the tree-decay exponent pays for that polynomial.
Raw II:11178--11195 uses this distinction between cubic and
merely vanishing contributions.

Changing the real cutoff nevertheless changes the integral.
In the original integration proof the Gaussian bad-bond gain
must be rechecked at the new threshold in (2.22), (2.26) and
(2.31). Its numerical scale is $e^{-cR(g)^2}$, which pays all
fixed geometric entropy below a sufficiently small coupling
ceiling. This numerical fact alone does not verify preservation
of the source's localized operator representations. The
boundary-specific domain CMP119 (3.43) is not a replacement
for the whole-lattice interior domain (1.34).

\subsection{Why the tunable density estimate matters}

The first overlap theorem above has a sufficient interaction norm
containing $K_G^{|Y|}$. If a source localization cube is one cell,
its real coordinate dimension is $D=O(M_{\rm loc}^4)$, and the
original $\log K_G$ has that dimension cost. Taking current-unit
cells instead produces the same loss through their number in
$Y$. The published coefficient decay does not automatically pay
this fixed cost: its parameter hierarchy may choose $M_{\rm loc}$
large after the decay exponent has been fixed. Decreasing the
coupling cannot cure an unfavorable exponential support exponent.

The tunable real-baseline comparison proved in this notebook
reduces that cost to $O(\theta D)$ per source cell. Choose
$\theta$ of order $D^{-1}$, with the source cubes as cells.
Its Cauchy derivative constant then grows with $\theta^{-1}$,
but is fixed before the coupling is decreased. This removes
the automatic $M_{\rm loc}^4$ loss in the support exponent.
The strict version also requires the displayed precision
perturbation smallness at that chosen $\theta$.

For example, if a source coefficient bound has sufficient spare
tree decay to pay the resulting fixed support weight and its
geometric summation, the Schwarz estimate above implies
$\epsilon_\theta(g)\le C_\theta gR(g)^3$.
The interaction and Gaussian-tail smallness conditions then hold
together below a fixed coupling ceiling. This is a precise
conversion criterion; the required source decay and precision
domain must still be matched. Alternatively one can retain
CMP116's own integration theorem, whose controllably small
quadratic and determinant errors serve the same purpose.

\subsection{The immediate model-specific proof}

The flat current-scale covariance is supplied by Notebook 17.
Notebook 19 constructs its nonflat physical-background version,
the exact relative-field pivot chart, and the local measure
factors on the homogeneous small-global-potential branch.
Notebook 20 verifies the full local coefficient and precision
match for the bare first-step Wilson small-field integral.
To invoke the original CMP116 theorem on its larger domains,
one must still show that the background, independent current
and every restricted reconstruction lie in its orbit domain,
and match the source's deformed operators.
If a cutoff or inherited carrier is changed, its exact transfer
through that proof must be established.

The retained cutoff coordinates in Notebook 19 are preserved
exactly. Additional indicators retain their own pullbacks.
The direct first-step proof does not establish the independent
current and mixed-history domains, transport of support to the
next coarse background, or the later quantum interactions.
Those are further model statements needed to close the induction.


\section{The model estimate to close next}
\label{oc18-next}

For a local analytic physical-field remainder $F_Y(z)$ vanishing
through degree two, bounded by $a_Y$ for $\norm z_{\infty,Y}<r_0$,
the one-variable Schwarz lemma on each complex radial line gives
\[
 |g^{-2}F_Y(gx)|\le a_Yr_0^{-3}gR^3
       \qquad (\norm x_{\infty,Y}\le R).
\]
If the corresponding weighted local sum of $a_Y$ is finite,
the interaction parameter in this notebook tends to zero with
$gR^3$. A merely extensive bound on the complete analytic action
does not establish that local sum. The full fluctuation exponent
also contains Jacobian terms beginning at first order and Haar
terms beginning at second order; their cutoff bounds must be
included with their actual prefactors.

Notebook 19 supplies the physical homogeneous Gaussian and its
local measure and gauge factors. Notebook 20 verifies the full
coefficient and precision hypotheses for the bare first Wilson
small-field integral. Notebook 21 supplies the localized
classical action and precision uniformly in earlier refinement
and closes one classical-reference integration. The remaining
model work is to control the inherited quantum interactions,
transport the activities to the next coarse background,
handle the complement of the small-field domain, and close the
renormalization induction. The continuum existence argument
requires that induction together with the observable and
reconstruction obligations of Notebook 01.

\begin{thebibliography}{9}
\bibitem{forest18}
A.~Abdesselam and V.~Rivasseau,
\emph{Trees, forests and jungles: a botanical garden for cluster
expansions}, \href{https://arxiv.org/abs/hep-th/9409094}{arXiv:hep-th/9409094}.
\bibitem{notes18}
A.~Abdesselam,
\emph{Notes on the Brydges--Kennedy--Abdesselam--Rivasseau
forest interpolation formula},
\href{https://mabdesselam.github.io/BKAR.pdf}{author notes},
Theorem 1 and the convex hull of partition masks.
\bibitem{kp18}
R.~Koteck\'y and D.~Preiss,
\emph{Cluster expansion for abstract polymer models},
Commun. Math. Phys. \textbf{103} (1986), 491--498,
\href{https://doi.org/10.1007/BF01211762}{doi:10.1007/BF01211762}.
\bibitem{prior18}
\emph{Toward a four-dimensional Yang--Mills construction},
research Notebooks 3--7 and 17 in this workspace.
The Gaussian cutoff and strict background-locality constructions
are separate inputs to the present extension.
\bibitem{balaban11618}
T.~Ba\l aban,
\emph{Renormalization group approach to lattice gauge field
theories. II. Cluster expansions},
Commun. Math. Phys. \textbf{116} (1988), 1--22,
\href{https://scispace.com/pdf/renormalization-group-approach-to-lattice-gauge-field-4penm4k7hc.pdf}{original article copy}.
\bibitem{balaban10918}
T.~Ba\l aban,
\emph{Renormalization group approach to lattice gauge field
theories I}, Commun. Math. Phys. \textbf{109} (1987), 249--301,
\href{https://scispace.com/pdf/renormalization-group-approach-to-lattice-gauge-field-3ygai0vxt1.pdf}{original article copy}.
\end{thebibliography}
\end{document}
